\chapter*{Conclusion}
\addcontentsline{toc}{chapter}{Conclusion}
In this thesis, we investigate the fractal structures in basin maps of the SCHW+BW and RN+MP systems. For the construction of the maps, we used the \textit{Gravitáček 2} solver in a slightly modified version and then applied our \textit{Python} scripts. For quantitative analysis, we calculated the fractal dimensions of the basin boundary using the uncertainty exponent method. 
First, we compared the results of our method with the results of \cite{Sychrovsky2020} and \cite{Polcar2019} for SCHW+BW. The comparison with \cite{Sychrovsky2020} was both qualitative and quantitative. We did a visual comparison of the basin maps. Then, we also calculated the fractal dimensions. Visually, our basins differed slightly. We could explain the difference by the fact that in our method the test particle can get as close to the ring as the relative tolerance allows. The difference was supported by quantitative analysis from the calculation of the fractal dimension. We also observed an effect of the integration time and precision. 
From this we can conclude that in analysis of dynamical systems, which are non-linear and also chaotic, even small details such as precision, integration method or integration time can significantly change the result. 
We also analyzed the dependence of the fractal dimension on the mass of the ring as was done in \cite{Sychrovsky2020}. The replication did not match exactly. Nevertheless, the trend was matched. This result showed us that our method can detect the main features, but in cases with very small features and high outcome sensitivity, a more careful analysis is needed, such as higher resolution or more perturbation values of $\varepsilon$.
Qualitative analysis was also applied in comparison with the Poincaré sections constructed in \cite{Polcar2019}. We observed a potential correspondence between some structures in the Poincaré sections and structures in the basin maps. That is, regular islands correspond with the "eternal orbit" basin, while white islands in the chaotic sea possibly correspond to islands in the "plunge/collision" basin. In addition, richer basins were observed in cases with more chaotic features in the Poincaré sections.
From these comparisons, we concluded that our method can provide meaningful results. Potentially, the biggest limitation of this method is the heuristics. A conservation-error termination is not evidence of a physical collision. The assumption of numerical stability may be true, as the validation suggests. However, stronger results may be obtained by closely investigating which trajectories are terminated due to the collision and which are terminated due to numerical failure, so they are potentially falsely classified. Other improvements can be made by increasing the resolution or increasing the number of perturbation values $\varepsilon$.
The RN+MP system was before analyzed, for example, with Poincaré sections or the Melnikov method, but not with basin map analysis. In this thesis, we brought this novelty.
In the RN+MP system, we observed the significant importance of precision. The accessible region of the RN+MP system tends to close towards the horizon. Therefore, we needed to analyze whether and when it is possible for a test particle to plunge into the black hole. Under carefully chosen parameters, the plunge into the black hole is possible without the accessible region being split into subregions. Motivated by \cite{PolcarSemerak2019}, we chose energy levels close to the separatrix and then increased them and analyzed the dependence of the fractal dimension on the energy. Both the basin maps and the decreasing fractal dimension with increasing energy are compatible with the homoclinic-chaos mechanism suggested by the analysis in \cite{PolcarSemerak2019}. It is therefore possible that transverse intersection of the stable and unstable manifolds is present, resulting in chaotic dynamics of homoclinic origin. To prove this hypothesis, further research is needed.
